\section{Ubicación del curso: qué es el aprendizaje por refuerzo profundo}

La primera clase de CS224R parte de la pregunta más fundamental: qué es el aprendizaje por refuerzo profundo (Deep Reinforcement Learning), en qué se distingue del aprendizaje automático tradicional y por qué vale la pena estudiarlo.

\begin{importantbox}{Definición del aprendizaje por refuerzo profundo}
El aprendizaje por refuerzo profundo estudia los \textbf{problemas de decisión secuencial} (Sequential Decision-Making Problems): el sistema necesita, con base en un flujo de información, alternar continuamente entre "observar" y "actuar". El curso estudia también las \textbf{soluciones} a este tipo de problemas, entre ellas imitation learning, online/offline RL, model-free/model-based RL, multi-task/meta RL, así como RL for LLMs y RL for robots. Todos los métodos insisten en poder escalar a redes neuronales profundas.
\end{importantbox}

\subsection{Diferencia central con el aprendizaje supervisado}

En el aprendizaje supervisado contamos con datos etiquetados $\{(x_i, y_i)\}$, y el objetivo es aprender una función $f(x) \approx y$. Los puntos de datos son independientes e idénticamente distribuidos (i.i.d.), y se nos indica \textbf{directamente} cuál es la salida correcta.

En el aprendizaje por refuerzo, la situación es completamente distinta:

\begin{itemize}
    \item \textbf{Objetivo}: aprender una política de comportamiento $\pi(a \mid s)$ que mapee estados a acciones.
    \item \textbf{Retroalimentación}: no se indica directamente la respuesta, sino que se aprende de la experiencia mediante \textbf{retroalimentación indirecta} (como una señal de recompensa).
    \item \textbf{Los datos no son i.i.d.}: la acción $a$ del agente afecta la observación futura $s'$, de modo que la distribución de los datos depende de la política actual.
\end{itemize}

\begin{knowledgebox}{Las múltiples formas del comportamiento}
El "comportamiento" en el aprendizaje por refuerzo profundo puede incluir: el control de movimiento de un robot, la conversación de un chatbot, la estrategia de un juego, las decisiones de un vehículo autónomo, las operaciones de un Web agent, entre otros. El rasgo común de estos escenarios es que la salida (acción) del modelo produce consecuencias en el entorno, lo que a su vez afecta las entradas futuras.
\end{knowledgebox}

\subsection{Alcance y objetivos del curso}

Chelsea Finn señala explícitamente que este curso no puede abarcar todo el contenido del aprendizaje por refuerzo profundo. El curso se centra en:
\begin{itemize}
    \item Los \textbf{conceptos centrales} detrás de los métodos de RL profundo, de modo que los estudiantes puedan entender métodos más avanzados.
    \item La \textbf{implementación} de los algoritmos, para que los estudiantes puedan programarlos por sí mismos.
    \item El \textbf{control de robots y los modelos de lenguaje} como ejemplos principales, aunque la técnica tiene una aplicabilidad mucho más amplia.
\end{itemize}

\subsection{Resumen de la sección}

El aprendizaje por refuerzo profundo es la disciplina que trata de la decisión secuencial. En comparación con el aprendizaje supervisado, el RL enfrenta retos particulares: la retroalimentación es indirecta, los datos no son i.i.d. y las acciones tienen consecuencias. El objetivo central del curso es que los estudiantes puedan entender e implementar los métodos de RL profundo, tanto los existentes como los emergentes.
